\section{O Axioma da União}
Por enquanto, dispomos apenas dos axiomas não lógicos da Extensão, da Compreensão e do Par.
Com esses axiomas, provamos a existência do conjunto vazio, conjuntos de um elemento e conjuntos de dois elementos.

Adiantamos, sem demonstração neste momento, o seguinte fato: apenas com esses axiomas, não é possível provar a existência de conjuntos com $3$ ou mais elementos.
Embora a demonstração formal desse fato só seja dada mais adiante, sua plausibilidade já pode ser percebida intuitivamente:
Com o Axioma do Par, garantimos a existência de conjuntos que contêm dois objetos previamente dados como elementos.
Com o Esquema da Compreensão, é possível formar, a partir de um conjunto já existente, um subconjunto definido por uma propriedade; ao menos aparentemente, isso não permite obter conjuntos com mais de dois elementos a partir apenas de conjuntos de até dois elementos.
Finalmente, o Axioma da Extensão não nos dá nenhuma ferramenta para construir novos conjuntos a partir de antigos.

Para provar a existência de conjuntos com mais de dois elementos, precisaremos de novos axiomas.
Um dos axiomas que nos permitirão ultrapassar essa limitação é o Axioma da União.
Tal axioma consta na formulação de 1908 de Zermelo \cite{zermelo1908untersuchungen} como seu Axioma V (\emph{Axiom der Vereinigung}).

\begin{axiom}[União]\index{Axioma da União}
    Para todo conjunto $x$, existe um conjunto $y$ que tem como elementos todos os elementos dos elementos de $x$.

    Em símbolos:

    \[
    \forall x \exists y \forall z\, (\exists w\, (w \in x \land z \in w)\rightarrow z \in y)
    \]
\end{axiom}

Esta formulação do Axioma da União garante a existência de um conjunto $y$ que tem como elementos todos os elementos dos elementos de $x$, mas não garante que $y$ tenha como elementos apenas os elementos dos elementos de $x$.
Contudo, em presença do Esquema da Compreensão e do Axioma da Extensão, obtém-se a caracterização exata da união, como enuncia o lema abaixo.
A demonstração é deixada como exercício (ver Exercício~\ref{ex:uniao}).

\begin{lemma}\index{Existência da União}\label{lem:existencia-uniao}
    Para todo conjunto $x$, existe um único conjunto $y$ tal que, para todo $z$, $z \in y$ se e somente se existe um $w$ tal que $w \in x$ e $z \in w$.

    Em símbolos:

    \[
    \forall x \exists! y \forall z\, (z \in y \leftrightarrow \exists w\, (w \in x \land z \in w)).
    \]
\end{lemma}

Devido ao lema acima, a definição abaixo se justifica.

\begin{definition}\index{união}
    Dado um conjunto $x$, a \emph{união de $x$} é o conjunto $\bigcup x$ definido como o único conjunto $y$ tal que, para todo $z$, $z \in y$ se e somente se existe um $w$ tal que $w \in x$ e $z \in w$.
    
    Formalmente, a união é regida pelo seguinte axioma definicional.
    \[
    \forall x \forall y\, \Bigl(y = \bigcup x \leftrightarrow \forall z\, (z \in y \leftrightarrow \exists w\, (w \in x \land z \in w))\Bigr).
    \]
\end{definition}

É possível que o leitor não esteja familiarizado com a notação de união de um conjunto, sendo esperado que o símbolo $\bigcup$ possua índices superiores ou inferiores.
Após definirmos formalmente famílias, definiremos formalmente tais notações.
Porém, pode ser útil para o entendimento do leitor adiantar que a notação $\bigcup x$ corresponde ao conjunto $\bigcup_{w \in x} w$.

A partir do já exposto, é possível definir a usual união entre dois conjuntos, $x\cup y$, como abaixo.

\begin{definition}
    Para quaisquer conjuntos $x$ e $y$, a \emph{união de $x$ e $y$}, denotada por $x\cup y$ é o único conjunto $z$ cujos elementos são exatamente os de $x$ e mais os de $y$.
    Formalmente, a união de dois conjuntos é regida pelo seguinte axioma definicional.

    \[
        \forall x \forall y \forall z\, \bigl(z = x \cup y \leftrightarrow \forall w\, (w \in z \leftrightarrow (w \in x \lor w \in y))\bigr).
    \]
\end{definition}

A definição acima requer justificativa, o que é feito no lema abaixo.

\begin{lemma}
    Para quaisquer conjuntos $x$ e $y$, existe um único conjunto $z$ cujos elementos são exatamente os de $x$ e mais os de $y$.

    Em símbolos:

    \[
        \forall x \forall y \exists! z\, \forall w\, (w \in z \leftrightarrow (w \in x \lor w \in y)).
    \]
\end{lemma}
\begin{proof}
    Para a existência, tome o conjunto $z = \bigcup \{ x, y \}$.
    Se $w \in z$, então existe um $w'$ tal que $w' \in \{ x, y \}$ e $w \in w'$.
    Assim, $w' = x$ ou $w' = y$, e então $w \in x$ ou $w \in y$.

    Por outro lado, se $w \in x$ ou $w \in y$, então $w \in w'$ para algum $w' \in \{ x, y \}$, e então $w \in z$.
    Assim, $w \in z$ se e somente se $w \in x$ ou $w \in y$.

    Para a unicidade, suponha que $z_1$ e $z_2$ sejam conjuntos tais que, para todo $w$, $w \in z_1$ se e somente se $w \in x$ ou $w \in y$, e $w \in z_2$ se e somente se $w \in x$ ou $w \in y$.
    Então, $\forall w \, (w \in z_1 \leftrightarrow (w \in x \lor w \in y) \leftrightarrow w \in z_2)$, e então $z_1 = z_2$ pelo Axioma da Extensão.
\end{proof}

Interessantemente, apesar de ser necessário a introdução de um novo axioma para garantir a existência de conjuntos com mais de dois elementos, o mesmo não ocorre para a definição de interseções.

Perceba que, para provar o lema abaixo, necessitamos apenas do Axioma da Extensão e do Esquema da Compreensão.
\begin{lemma}
    Seja $x$ um conjunto não vazio.
    Então, existe um único conjunto $y$ tal que, para todo $z$, $z \in y$ se e somente se $z$ é um elemento de todos os elementos de $x$.

    Em símbolos,

    \[
        \forall x\, (x \neq \emptyset \rightarrow \exists! y\, \forall z\, (z \in y \leftrightarrow \forall w\, (w \in x \rightarrow z \in w))).
    \]
\end{lemma}
\begin{proof}
    Fixado $x\neq \emptyset$, tome um elemento $w'$ de $x$.
    Agora, tome o conjunto $y = \{ z \in w' : \forall w\, (w \in x \rightarrow z \in w) \}$.
    Como $w' \in x$, é imediato que para todo $z$, $z \in y \leftrightarrow \forall w\, (w \in x \rightarrow z \in w)$.

    A unicidade é uma consequência imediata do Axioma da Extensão.
\end{proof}

\begin{definition}\index{interseção}\label{def:intersecao}
    Seja $x$ um conjunto não vazio.
    A \emph{interseção de $x$} é o conjunto $\bigcap x$ definido como o único conjunto $y$ cujos elementos são exatamente os elementos comuns a todos os elementos de $x$.

    Formalmente, a interseção obedece o seguinte axioma definicional.

    \[
        \forall x\, (x \neq \emptyset \rightarrow \forall y\, (y = \bigcap x \leftrightarrow \forall z\, (z \in y \leftrightarrow \forall w\, (w \in x \rightarrow z \in w)))).
    \]
\end{definition}

Uma observação fundacional sem impactos práticos para o desenvolvimento da teoria, mas necessária, é a de que, de acordo com a Metadefinição~\ref{def:extensao-por-definicao}, para se introduzir por definição um símbolo funcional $n$-ário a uma teoria $T$, é necessário que $T\Vdash\forall x_1 \ldots \forall x_n \exists! y\, \varphi(x_1, \ldots, x_n, y)$, onde $\varphi$ é uma fórmula da linguagem de $T$.
Feito isso, insere-se um símbolo funcional $n$-ário $f$ e o seguinte axioma definicional:
\[
\forall x_1 \ldots \forall x_n \forall y\, (y = f(x_1, \ldots, x_n) \leftrightarrow \varphi(x_1, \ldots, x_n, y)).
\]

No caso, $\bigcap$ é um símbolo funcional unário.
Seja $\varphi(x, y)$ a fórmula $\forall z\, (z \in y \leftrightarrow \forall w\, (w \in x \rightarrow z \in w))$.
Note que nós \emph{não} provamos que $\forall x \exists! y\, \varphi(x, y)$, mas apenas que $\forall x\, (x \neq \emptyset \rightarrow \exists! y\, \varphi(x, y))$.
Em verdade, para $x=\emptyset$, tal afirmação é falsa (ver Exercício~\ref{ex:uniao-sem-intersecao}).

Para que a introdução do símbolo funcional $\bigcap$ seja legítima dentro do nosso formalismo, é necessário atribuir algo arbitrário para $\bigcap \emptyset$, ainda que não valha que $\bigcap \emptyset$ seja o conjunto cujos elementos são exatamente os elementos comuns a todos os elementos de $\emptyset$.

Essa ocorrência está longe de ser exclusiva deste caso em particular.
Por exemplo, mais adiante, definiremos a noção de número ordinal --- uma generalização da noção de número natural.
Conforme veremos, será definida uma soma para ordinais a partir de uma introdução por definição de um símbolo funcional binário $+$.
Apesar de só ser desejável definir $\alpha+\beta$ quando $\alpha$ e $\beta$ são ordinais, a introdução por definição de $+$ exigirá que se atribua algo arbitrário para $\alpha+\beta$ quando $\alpha$ ou $\beta$ não sejam ordinais.

Uma solução padrão para este tipo de problema é definir, de forma arbitrária, um valor qualquer, porém fixo, para os casos em que se gostaria de deixar o símbolo funcional ``indefinido''.

Para tanto, reformularemos, abaixo, a definição de interseção, atribuindo um valor arbitrário para $\bigcap \emptyset$.

\begin{definition}\index{interseção}\label{def:intersecao2}
    Inserimos o símbolo funcional unário $\bigcap$ por definição, a partir do seguinte axioma definicional.

    \[
        \forall x\, \forall y\, \bigl(y = \bigcap x \leftrightarrow ((x = \emptyset \land y = \emptyset) \lor (x \neq \emptyset \land\forall z\, (z \in y \leftrightarrow   \forall w\, (w \in x \rightarrow z \in w))))\bigr).
    \]
\end{definition}

No caso, atribuímos o valor arbitrário de $\emptyset$ para $\bigcap \emptyset$.
Tal atribuição não possui consequências na teoria, uma vez que sempre evitaremos intersectar o conjunto vazio (o que faríamos de qualquer modo caso ele estivesse ``indefinido'').
É importante ressaltar que $\bigcap x$ apenas é o conjunto cujos elementos são exatamente os elementos comuns a todos os elementos de $x$ quando $x$ é um conjunto não vazio.
\section*{Exercícios}
\begin{exercise}\label{ex:uniao}
    Prove o Lema~\ref{lem:existencia-uniao}.
\end{exercise}

\begin{exercise}
    Seja $X$ um conjunto.
    Prove que $\bigcup X$ corresponde ao supremo de $X$ na ordem delimitada por $\subseteq$ no universo dos conjuntos.
    Ou seja, prove os itens abaixo.
    \begin{enumerate}
        \item Para todo $x \in X$, $x \subseteq \bigcup X$.
        \item Se $Y$ é um conjunto tal que, para todo $x \in X$, $x \subseteq Y$, então $\bigcup X \subseteq Y$.
    \end{enumerate}
\end{exercise}

\begin{exercise}\label{ex:uniao-sem-intersecao}
    Prove que, não existe $y$ tal que, para todo $z$, $z \in y$ se e somente se $\forall w\, (w \in \emptyset \rightarrow z \in w)$.
    Mais informalmente, mostre que $\{x: \forall w\, (w \in \emptyset \rightarrow x \in w)\}$ não é um conjunto.
\end{exercise}